\documentclass{article}
\usepackage{pathfinder-note}

\title{Priced Certificates for Conditional-CVaR Obedience}

\begin{document}
\maketitle

\section{The pair}

Q surveys learning-augmented algorithms and asks how prediction cost, feedback, and component guarantees should be accounted for in a system. P studies fixed-batch CVaR values in a one-type obedience test, showing that expected empirical CVaR can reverse an exact state-independent reward decision; the mechanism-design paper called Q inside P is a different paper from the supplied Q. \ledger{1, 4, 5}

\section{The connexion pursued}

The peers asked whether a \emph{realised}, priced batch of conditional-loss samples could certify P's obedience decision, then whether that certificate could enter Q's cost and composition framework. The narrow setting is a fixed policy on finite signals, with utility equal to negative conditional upper-tail CVaR, state-independent action rewards, and access to losses for prescribed actions and deviations. It does not cover ex-ante CVaR of a plan, report-stage incentives, or a general learning-augmented pipeline. \ledger{3, 4, 5, 15, 20}

P already supplies two ingredients: a reversal for every finite batch when decisions use \emph{expected} empirical CVaR, and deterministic cycle and reward tests given a bound on every conditional utility. Its final limitation leaves certification from one realised batch unresolved. Q supplies accounting questions, including a charge per predictor invocation and conditions for composing component guarantees, but no obedience theorem. \ledger{4, 5, 8, 20}

\section{What was found}

\begin{claim}
A realised batch can give a checkable, high-confidence certificate of exact true obedience, provided the relevant conditional losses can be sampled. A failed certificate can require abstention.
\end{claim}

Let \(S\) be the signal set, \(B\) the available prescribed actions, \(N=|S||B|\), and \(a(s)\in B\) the policy. For each cell \((s,b)\), draw \(m\) independent losses from a fixed distribution on \([0,H]\). Set
\[
U(b,s)=-\operatorname{CVaR}_{\beta}(L_{s,b}),
\qquad
\widehat U(b,s)
=-\operatorname{CVaR}_{\beta}(\widehat F_{s,b}),
\]
where \(\beta\in(0,1]\) is upper-tail mass and \(\widehat F_{s,b}\) is the realised empirical distribution. The CVaR variational formula bounds risk error by \(H\|F-\widehat F\|_{\infty}/\beta\). Dvoretzky--Kiefer--Wolfowitz and a union bound therefore give simultaneous utility error at most
\[
e_m=\frac{H}{\beta}
\sqrt{\frac{\log(2N/\delta)}{2m}}
\]
with probability at least \(1-\delta\). Independence across cells is unnecessary for this union bound. This interval concerns realised data; P's expected empirical value is a different quantity. \ledger{3, 8, 20}

For distinct prescribed actions \(a,b\), compute
\[
w_m(a,b)=
\max_{s:\,a(s)=a}
\{\widehat U(b,s)-\widehat U(a,s)+2e_m\}.
\]
If this directed graph has no positive cycle, its difference constraints have a solution
\[
r(a)-r(b)\geq w_m(a,b).
\]
The recorded estimates, radius, graph check, and reward vector form a checkable certificate. On the simultaneous-coverage event, the same \(r\) satisfies every true obedience inequality. Conversely, a displayed length-\(K\) prescribed-action cycle with empirical utility sum \(\widehat C\) and \(\widehat C+2Ke_m<0\) certifies true infeasibility on that event. If neither test succeeds, the answer is abstention. The graph step instantiates P's deterministic rectangular-error test; it is not a new graph theorem. \ledger{8, 14, 20}

If every nontrivial true simple cycle has utility sum at least \(K\gamma\), the feasible certificate succeeds on the coverage event when \(e_m<\gamma/4\). A true cycle with sum at most \(-K\gamma\) likewise yields a negative certificate under that condition. Thus uniform sampling with
\[
m>
\frac{8H^2}{\beta^2\gamma^2}
\log\frac{2N}{\delta}
\]
is sufficient for either stated margin promise. If each loss draw costs \(\kappa_s\), the acquisition bill is \(\kappa_sNm\), plus any material graph-computation charge. Equating a draw with Q's priced predictor invocation is an additional model choice. \ledger{8, 15, 17, 18}

\begin{claim}
A reward witness can be compared with the cheapest true-utility reward cap, but sampling cost alone does not bound implementation cost.
\end{claim}

Let \(w^*(a,b)\) be the true counterpart of \(w_m(a,b)\), and suppose every nontrivial simple cycle has average true edge weight at most \(-\gamma\). On coverage,
\[
w^*(a,b)\leq w_m(a,b)\leq w^*(a,b)+4e_m.
\]
For nonnegative rewards, the coordinatewise least solution is the maximum weight of a path starting at each action, including the empty path. With no positive cycle, a maximizing path can be simple and has at most \(n-1\) edges, where \(n=|B|\). If \(R^*\) is the smallest true-utility maximum reward and \(R_m^*\) the smallest confidence-robust maximum reward, then
\[
R^*\leq R_m^*
\leq R^*+4(n-1)e_m
\qquad\text{when }e_m<\gamma/4.
\]
When one reward payment counts as resource cost, this gives the confidence-qualified local bound
\[
\kappa_sNm+R^*+4(n-1)e_m.
\]
It is a comparison with a local support-payment oracle, not a competitive ratio against Q's offline system benchmark. \ledger{10, 11, 12, 14, 18}

The distinction matters even without sampling: constraints
\[
r(a)-r(b)\geq M,
\qquad
r(b)-r(a)\geq-M-\gamma
\]
have a nonpositive cycle, yet any nonnegative supporting reward has maximum at least \(M\). A declared reward budget can instead be checked by adding \(0\leq r(a)\leq R\) to the confidence-robust difference constraints. Failure of that robust budget check alone does not prove that the true array violates the budget. \ledger{10, 11, 12}

\section{Objections and their status}

\begin{objection}
Without a margin, finite samples cannot uniformly decide exact feasibility; without deviation observations, even unlimited on-policy data may not identify it.
\end{objection}

For the finite-sample obstruction, use two swapped-action signals, safe loss \(1\), and risky loss \(2X\) with \(X\sim\operatorname{Bernoulli}(p)\). At \(\beta=1/2\) and \(p<1/2\), risky CVaR is \(4p\), so the true two-edge cycle is \(2-8p\). It is positive at \(p_-=1/4-\varepsilon\) and negative at \(p_+=1/4+\varepsilon\). Every finite binary transcript has positive probability in both worlds. For a nonabstaining test with error at most \(\delta\) in each, the transcript change-of-measure bound gives
\[
\mathbb E_{p_-}T\,
\operatorname{kl}(p_-,p_+)
\geq
\operatorname{kl}(1-\delta,\delta),
\]
where \(T\) counts informative risky draws. Since the per-draw divergence is of order \(\varepsilon^2\), the required expected sampling bill grows as the cycle margin tends to zero. Abstention and a positive-margin promise remain available. \ledger{6, 7, 15, 20}

For the identification obstruction, make each prescribed loss constantly \(3/2\). Let every unplayed deviation cost \(2\) in one world and \(1\) in another. The on-policy logs have the same law at every sample count, while the true two-edge utility cycles are \(+1\) and \(-1\). The certificate therefore needs a generative conditional sampler, covered exploration, or a justified structural identification model; its cost must be included. \ledger{9, 13, 15, 20}

\begin{objection}
A \(1-\delta\) certificate is not an expected-cost or pipeline guarantee.
\end{objection}

An expected-cost statement needs a specified abstention action and a bound on the coverage-failure event. If downstream cost \(D\) is at most \(B\) on the \emph{entire} coverage event, including abstention, and at most \(M\) on every path, then
\[
\mathbb E[\kappa_sNm+D]
\leq\kappa_sNm+B+\delta M.
\]
The peers corrected an earlier formulation that bounded only accepted certificates. In Q's adaptive pipeline, coverage and cost premises must hold conditionally on every reachable pre-sampling history; a marginal statement for one fixed policy is insufficient. A competitive ratio further requires a common positive benchmark and relations between component instances and system cost. Those relations were not established. \ledger{20, 22, 24}

\section{Outside sources and what remains open}

The peers used existing empirical-CVaR concentration and risk-arm identification work, including \url{https://arxiv.org/abs/1901.00997}, \url{https://arxiv.org/abs/1902.10709}, and \url{https://proceedings.neurips.cc/paper/2021/file/d69c7ebb6a253532b266151eac6591af-Paper.pdf}. They also cited a standard sequential change-of-measure source at \url{https://arxiv.org/abs/1407.4443} and related incentive-margin work at \url{https://arxiv.org/html/2605.12094v1}. These sources limit a novelty claim for concentration plus P's graph test; the recorded searches did not establish priority for the combined example. \ledger{7, 13, 20}

The pair supports a narrow answer: priced counterfactual samples can yield a high-confidence true-obedience witness and a local reward-cap bound under stated access and margin assumptions. The material is thin on any broader Q-style system guarantee. Open points are a justified counterfactual access model, tighter graph-aware adaptive sampling costs, treatment of transfers in the objective, and a bounded abstention or failure path tied to a common offline benchmark. \ledger{15, 20, 24}

The smallest next step toward the system claim is to specify one downstream objective, its prediction-free benchmark, whether rewards count as costs, and a feasible abstention fallback in the two-signal model. Proving the resulting conditional expected-cost comparison, or giving a counterexample, would settle whether this certificate composes in that minimal setting. \ledger{15, 20, 22, 24}

\end{document}